The module is built out of tree against the headers of the running kernel:
the Makefile hands \texttt{M=\$(CURDIR)} to the kernel's own build system in
\texttt{/lib/modules/\$(uname -r)/build}. The compiler warning compares two
names for the same GCC~13.3.0 and is harmless. BTF generation is skipped
because the headers package carries no \texttt{vmlinux}; that affects only
the type information used by BPF tools. The module is not signed and is built
outside the kernel tree; the kernel loads it because signatures are not
enforced, but marks itself tainted. \texttt{/proc/sys/kernel/tainted} reads
\Tainted{}: 4096 for an out-of-tree module and 8192 for an unsigned one, both
from \texttt{pfcount}, plus 16384 for the soft lockups the kernel reported
while booting, before the module was loaded. Those come from the slow
virtualisation of this VM (Mini-Project~1), not from the module. The test programs
compile without warnings under \texttt{-Wall -Wextra}.
